% Source: https://artofproblemsolving.com/community/c3h2536374_quintessential_quintuplets
% Stable identity: aops-topic-2536374-q01
The conditions give $32\ge t\ge16i$, so $i=1$ or $i=2$.

\emph{$i=2$:} then $n\ge4$, $m\ge8$, $y\ge16$, $t\ge32$, which forces $(2,4,8,16,32)$: $1$ quintuplet.

\emph{$i=1$:} we need $2\le n$, $2n\le m$, $2m\le y$, $2y\le t\le32$, which forces $y\le16$, $m\le8$, $n\le4$. For fixed $y$ there are $33-2y$ choices of $t$. For fixed $m$, summing over $y=2m,\dots,16$ adds the odd numbers $33-4m,\dots,3,1$, giving $(17-2m)^2$. So $m=4,5,6,7,8$ contribute $81,49,25,9,1$, and
\[n=2:\ 81+49+25+9+1=165,\qquad n=3:\ 25+9+1=35,\qquad n=4:\ 1,\]
for a total of $201$.

Altogether there are $201+1=\boxed{\textbf{(B) } 202}$ quintessential quintuplets.
